\scanpage{43}
% archive sheet number: 42

\textcircled{\scriptsize 12} Anneau compact (i.e. à spectre compact)\quad
$X_{\rm cpct}\simeq\widetilde R_{\mathrm{top\,cpct}}$

$\forall U$, $A(U)$ à spectre cpct i.e.
$\forall f\in A(U)$, $\exists(1)\ e\in A(U)$, $e^2=e$, $ef$ inversible dans $eA(U)$,
\deleted{$e$} $(1-e)f$ nilpotent. Suffit de le faire pour $f$ universel, $U=A$.
On trouve, si $D_n\subset A\times A\times A$ formé des $g,f,e$ avec $e^2=e$
$f^n(1-e)=0$, $g(1-e)=0$, $fg=e$, \unsure{on voit|on voit se|un mono} induit par \deleted{\unsure{pr}} \unsure{proj$_2$}

\[
 i_n:D_n\hookrightarrow A
\]

et la condition est que les $D_n$ couvrent $A$

Pour la top du $R_{\rm top\,cpct}$, sert donc la condition que, si
$D_n=I_n\amalg_R A^*\lhook\joinrel\xrightarrow{i_n}A$
\inserted{attention, c'est une somme \uline{ds $R$}, pas $\widehat R$}
la famille des $i_n$ soit couvrante.

\margin{topologie non cohérente ! A-t-elle assez de pts ??}\quad ?

\textcircled{\scriptsize 13} \uline{Anneau compact réduit} $=$ \uline{pseudo\deleted{\unsure{cor|corps}}corps}
\quad $X_{\rm cpct}\simeq\widetilde R_{\mathrm{top\,ps\,corps}}\simeq\widehat R_{\rm cons}$

$\forall U$, $A(U)$ est ps. cpct i.e. $\forall f\in A(U)$, $A(U)\xrightarrow{\sim}A_f\times A/fA$, i.e.
$\exists(1)\ e\in A(U)$, $e^2=e$, $ef$ inv., $e(1-f)=0$
i.e. $i_1:D_1\hookrightarrow A$ est un isom

ps. corps est la topologie engendrée par les $X'\to X$ qui ont une section fibre par fibre.
(($X_i\to X$) couvrante ssi $\exists\deleted{\unsure{i|i_0}} X_{i_0}\to X$ ainsi)

$R_{\rm cons}$ comme avant, localisée de $R$ pour les \deleted{topologie} flèches \deleted{\unsure{bijectives|bijections}}
à ext. résiduelles triviales.

$\widehat R_{\rm cons}\subset\widehat R_{\unsure{\rm cons|\rm ét|R|\rm cpct}}$ formé des $R^\circ\to\mathrm{Ens}$ qui transforment les flèches précédentes

\textcircled{\scriptsize 14} Corps sép. clos $=$ corps \inserted{str.} hensélien
\quad $X_{\rm corps\,hens}\simeq\widetilde R_{\mathrm{top\,corps\,sépclos}}$

top. corps hens, topologie où les familles couvrantes sont celles \deleted{\unsure{pour lesquelles}}
$X_i\to X$ qui \deleted{\unsure{sont de fibre séparablement}} ont un pt \deleted{\unsure{rationnel}}
séparable (\unsure{ou fini sép.|ou pure$^t$ sép.|i.e. (fini sép.)}) sur chaque $x\in X$

\[
 \simeq\widetilde{(R_{\rm cons})}_{\rm et}\qquad R_{\rm cons}\text{ comme dans 11), avec top. \uline{étale}.}
\]
